\section{Guaranteed success does not elicit private costs}
Consider $n$ agents with independent positive finite-support private costs
$C_i<B$. An active set $S$ succeeds exactly when $S\in\mathcal F$, where
$\mathcal F$ is a known nonempty upward-closed family and
$\varnothing\notin\mathcal F$. Success gives every agent benefit $B$;
agent $i$ pays $C_i$ iff active. There are no transfers, verification or
type-dependent action restrictions. A finite communication protocol may use
public and private randomization independent of costs. We consider an equilibrium
of the whole reporting, communication and action game, not obedience to a
mediator who already knows the costs.

\begin{proposition}[Guaranteed-success frontier]
Among such equilibria succeeding surely at every supported type profile,
maximum expected welfare is
\[
 nB-\min_{S\in\mathcal F}\sum_{i\in S}\mathbb E C_i.
\]
First-best welfare is attainable exactly when some fixed $S\in\mathcal F$
is realized-cost-minimizing almost surely.
\end{proposition}
\begin{proof}
Let $x_i(c)$ be agent $i$'s interim probability of activity. A type $c$ can
mimic another type $c'$'s entire randomized communication and action strategy.
Independence makes the opponents' type law the same, and guaranteed success
makes this deviation's success benefit $B$. Hence
$B-cx_i(c)\ge B-cx_i(c')$. Since $c>0$, interchanging types gives
$x_i(c)=x_i(c')=:x_i$. Thus expected cost is
$\sum_i\mathbb E[C_i]x_i$. The unconditional active-set law is supported on
$\mathcal F$, so this sum is at least the minimum displayed coalition cost.
Choose a coalition attaining that minimum. Positive means imply it is
inclusion-minimal: its members are pivotal, prefer $B-C_i>0$ to failure,
and outsiders prefer abstention. This designated-coalition equilibrium attains
the bound. Finally the minimum of expected coalition costs equals the expected
minimum realized cost iff a fixed minimizing coalition attains that realized
minimum almost surely, since their nonnegative difference must vanish.
\end{proof}

This is a finite-protocol specialization of established screening reasoning,
not a new general incentive principle. Battaglini--Palfrey's honest-and-obedient
constraints already include joint report/action deviations\citep{battaglini2024}.
The statement covers threshold and coverage tasks but requires independence,
guaranteed success and the stated common-benefit utility. It is weak
implementation of one equilibrium, not uniqueness or full implementation.

For a calibration, let $n=2$, $C_i\in\{l,h\}$ independently, with low
probability $p$, $0<l<h<B$, $a=1-p$ and $\bar c=pl+ah$.
The unconstrained team optimum over labelled local type policies and
type-independent public randomness is
\[
 W_{\rm PL}=\max\{2B-\bar c,\;2B(2p-p^2)-2pl\}.
\]
Indeed a multilinear objective attains its maximum at deterministic local
policies. Conditional on the other's policy, activity is worthwhile at cost
$c$ iff $c\le B\Pr(\text{other inactive})$; thus the high-only rule is
unnecessary, and checking the remaining vertices gives the formula.
An externally informed obedient mediator attains
$W_*=2B-l(2p-p^2)-ha^2$, with gap
\[
 W_*-W_{\rm PL}
 =\min\{pa(h-l),\;(2B-h)a^2+lp^2\}>0.
\]
Equality occurs only at $p\in\{0,1\}$ or $h=l$ in this closed parameter
domain. Neither this team bound nor the externally informed mediator is a
truthful cheap-talk implementation. With transfers conditional on verified
completion, second-price procurement admits an ex-post honest-and-obedient
equilibrium; arbitrary disobedient opponents prevent the stronger claim of
dominant strategies in the full voluntary-action game.

\section{Pricing trusted verification}
Now restrict the architecture: the coordinator receives no private reports or
other type-dependent signals, and can query an external truthful verifier for
one previously uninspected actual cost per unit. Histories, stopping and
selection depend only on those answers and independent randomness. After
inspection it designates one worker, who obeys since $B>C_i$.
For arbitrary $n\ge2$ iid binary costs as above, put
\[
 W_0=nB-\bar c,\quad v=pa(h-l),\quad
 R_*=\sum_{j=0}^{n-2}a^j=\frac{1-a^{n-1}}p.
\]
\begin{proposition}[Exact expected-query frontier]
For expected query budget $r\ge0$, maximum gross welfare is
$W(r)=W_0+v\min\{r,R_*\}$.
With a hard cap $m\le n-1$, it is $nB-l-(h-l)a^{m+1}$.
Zero-error first best requires worst-case $n-1$ queries and minimum expected
count $R_*$.
\end{proposition}
\begin{proof}
Before finding a low cost, with at least two uninspected agents, the best
fallback has expected cost $\bar c$. A query finds low with probability $p$
and reduces cost to $l$; otherwise an uninspected fallback remains. Each such
useful query has conditional gain $p(\bar c-l)=v$. After finding low, or with
only one uninspected agent, queries have zero value. If $U$ counts useful
queries, conditional expectation gives gain at most $v\mathbb E U$.
The $j$th useful query occurs with probability at most $a^{j-1}$, so
$\mathbb E U\le\min\{r,R_*\}$. Search until the first low or $n-1$ highs,
then select the low or uninspected fallback. This attains first best with
expected count $R_*$. Mixing this whole policy with no search attains every
intermediate expected-budget point. Capping search at $m$ gives expected
count $\sum_{j=0}^{m-1}a^j$ and the stated value. Finally, after fewer than
$n-1$ high observations, stopping can select a high while an uninspected low
exists with positive probability. Zero-error selection must continue on all
such histories, proving both lower bounds.
\end{proof}
A uniform random inspection order gives an equivariant law with unchanged
query count; full search selects uniformly among cheapest agents.
For per-query price $\kappa\ge0$, net optimum is
$W_0+\max\{0,v-\kappa\}R_*$. Query addresses, role labels and recommendation
bits are outside this resource count. This is a special nonobligatory-inspection
problem\citep{beyhaghi2022}, not a general minimum-communication theorem.

\section{Action-orbit algorithms and the input boundary}
Consider complete-information rational games with explicitly listed actions
and payoffs $u_i(s)=\sum_{F\in\mathcal F_i}f_{iF}(s_F)$, where $i\in F$,
$|F|\le d$ and each factor table is explicit. Let $L$ include all table,
scope, coefficient and supplied typed-group generator encodings. Exact
automorphism validity is a promise; its discovery is not free.
Let $q$ count group orbits of the disjoint player-action pairs $(i,a)$.
An invariant independent profile assigns probability $x_w$ to each action in
orbit $w$, with normalizations $\sum_w c_{iw}x_w=1$.

\begin{proposition}[Local-factor orbit reduction]
Conditional payoff polynomials $P_w(x)$ of degree at most $d-1$ can be built
in polynomial time in $L$. Invariant Nash equilibria are exactly solutions of
\[
 x_w\ge0,\quad \sum_w c_{iw}x_w=1,\quad
 P_w(x)\le v_{B(w)},\quad
 x_w(v_{B(w)}-P_w(x))=0,
\]
where $B(w)$ is the player orbit. There are at most $2q$ variables and
$O(q)$ constraints of degree at most $D=\max\{2,d\}$.
\end{proposition}
\begin{proof}
For a representative $(i,a)$, sum each factor table against the independent
probabilities of its other players, fixing $a$. This constructs $P_w$ with
at most $d-1$ probability factors per term. Exact typed automorphisms identify
conditional payoffs throughout the orbit. There are at most $q$ player orbits;
the displayed constraints say precisely that every action's payoff is at most
the player's value and every positive-probability action attains it. These
are the mixed Nash conditions. The classical invariant-equilibrium existence
argument applies, so the system is nonempty.
\end{proof}

Using standard real-algebraic sampling, this gives runtime
$(qD)^{O(q)}\operatorname{poly}(L)$ for an algebraically represented exact
solution. The coefficient-bit and output-encoding bound is an external
algorithmic ingredient, not implemented here; its precise original-source
formulation remains to be verified before publication. The primary algorithmic
baseline is Tewolde et al.'s action-orbit method\citep{tewolde2025}; symmetric
graphical-game algorithms are also prior art\citep{brandt2011}.

There is a useful input-size refinement. Remove singleton-action players and
substitute their actions. A residual factor on $t$ active players has at least
$2^t$ table entries, so $t\le\log_2 L$. Consequently the bound is
$(q\log L)^{O(q)}\operatorname{poly}(L)=f(q)\operatorname{poly}(L)$:
for $a=Cq$, $(\log L)^a\le(a/e)^aL$, up to fixed log-base factors.
The same observation applies to the explicit normal-form baseline; it does
not establish a new separation from that prior algorithm. Expanding the
orbit coordinates to individual actions costs at least the output length.
Rational approximation needs a normalization repair: approximate orbit
masses within each player's action-orbit blocks, renormalize their sum, and
divide by block multiplicities. Payoff Lipschitz bounds from the explicit
rational tables give an $\varepsilon$-Nash output at precision polynomial in
$L$ and $\log(1/\varepsilon)$.

One can also append a polynomial welfare threshold to this compact system
and bisect its feasible values, obtaining an $\eta$-optimal equilibrium
\emph{within the invariant Nash set}. Compactness gives attainment; rational
rounding requires an additional welfare-error allowance. This does not
optimize unrestricted Nash or guarantee good welfare.

An invariant correlated-equilibrium program instead uses probabilities on
\emph{profile} orbits and representative deviation-orbit inequalities.
Averaging preserves CE and an invariant linear welfare objective; an optimal
solution can use at most $r+1$ profile orbits when there are $r$ deviation rows.
However binary cyclic games can have at least $2^n/n$ profile orbits despite
only two action orbits. Neither enumeration nor decentralized correlated
sampling is free. Compact CE algorithms are established prior
work\citep{papadimitriou2008,jiang2011}.

\subsection{A one-orbit counting obstruction}
Given a Boolean circuit $F$ on $n\ge2$ bits, let $r$ rotate coordinates and
$c$ complement them. Give all players the same succinct payoff circuit
\[
 u_i(s)=\frac1{2n}\sum_{j=0}^{n-1}\bigl(F(r^js)+F(cr^js)\bigr).
\]
Rotation and simultaneous complementation are exact typed automorphisms,
transitive on player-action pairs: $q=1$. The unique invariant product profile
is fair independent bits. It is Nash because complementation equates each
player's two conditional payoffs. Yet
\[
 \mathbb E u_i=2^{-n}\#\{s:F(s)=1\}.
\]
Every transformation permutes the input cube, proving the equality and an
explicit counting-SAT reduction for exact value evaluation. This standard-style
reduction concerns succinct circuits, not explicit local tables. Finding the
uniform equilibrium is trivial here; invariant regret is zero. No hardness
of equilibrium output, regret evaluation or value-free near-optimal profile
selection follows.

\section{When symmetry averaging changes implementation}
Let a finite mechanism's feasible vector set be nonempty compact convex and
preserved by a finite relabelling group. Include linear joint report/action
incentives and a group-invariant linear resource budget. For an invariant
\emph{affine} objective $J$, averaging an optimizer over the group gives an
invariant optimizer at the same resource level. This is established reasoning,
explicit in collective-action mechanism design\citep{battaglini2024}.
Invariance alone is insufficient: on the two-coordinate probability simplex,
$J(z)=z_1^2+z_2^2$ has only noninvariant vertex maximizers under exchange.

Nor does averaging preserve product-law implementations. Independent actions
produce exactly one active agent with probability one only by designating a
surely active agent and making all others surely inactive. Their permutation
average is correlated, so public role selection changes the architecture.
At a fixed input, allowed outputs decompose into stabilizer orbits. An invariant
law is a mixture of uniform orbit laws, with minimum entropy $\log_2 m$ for
smallest orbit size $m$. A deterministic decoder of $b$ fair seed bits exists
exactly when $2^b$ is a nonnegative integer combination of allowed orbit sizes.
A single three-element orbit therefore has no fixed finite fair-bit sampler;
rejection using two bits has expected cost $8/3$ bits and unbounded worst case.
These are centralized sampling facts, not general distributed communication
or anonymous-election bounds.
